\section{Integrales en una dimensión}
\label{sec:mc-integrales}

Se estiman las dos integrales del enunciado,
\[
  \text{(a)}\;\int_0^1 \sin(\pi x)\,dx = \frac{2}{\pi} \approx 0.63661977,
  \qquad
  \text{(b)}\;\int_0^2 \frac{1}{\sqrt{2\pi}}e^{-x^2/2}\,dx
  = \Phi(2)-\Phi(0) \approx 0.47724987,
\]
usando el estimador \eqref{eq:estimador} y el intervalo \eqref{eq:ic}. Cada
integral se calcula con dos fuentes de uniformes: el MT19937 implementado en la
Parte~\ref{sec:fundamento} y el LCG de Numerical Recipes.

\begin{table}[H]
\centering\small
\begin{tabular}{llrrcrrc}
\toprule
Integral & Fuente & $\hat I_N$ & EE & IC 95\% & Valor real & Error & Contiene \\
\midrule
(a) & MT19937 (propio) & 0.637372 & 0.000973 & [0.635466, 0.639279] & 0.636620 & +7.53e-04 & Sí \\
(a) & LCG (propio) & 0.637499 & 0.000974 & [0.635590, 0.639407] & 0.636620 & +8.79e-04 & Sí \\
(b) & MT19937 (propio) & 0.477405 & 0.000727 & [0.475980, 0.478831] & 0.477250 & +1.55e-04 & Sí \\
(b) & LCG (propio) & 0.476433 & 0.000727 & [0.475008, 0.477858] & 0.477250 & -8.17e-04 & Sí \\
\bottomrule
\end{tabular}

\caption{Estimaciones con $N = 100\,000$. Los cuatro intervalos contienen el
valor verdadero.}
\label{tab:mc-integrales}
\end{table}

Los cuatro intervalos atrapan el valor real. Los errores son del orden de
$10^{-4}$, consistentes con un error estándar de $\sim 10^{-3}$: las
estimaciones caen a menos de un error estándar del valor exacto, que es
exactamente lo que predice la teoría.

\subsection{Estudio de convergencia}

La Figura~\ref{fig:convergencia} grafica el error absoluto medio contra $N$ en
escala log--log, con $N$ desde $10^1$ hasta $10^6$. Para que la recta sea
legible se promedia el error sobre 20 réplicas en bloques disjuntos: una sola
corrida produce una nube demasiado ruidosa para ajustar una pendiente con
sentido.

\begin{figure}[H]
\centering
\includegraphics[width=\textwidth]{convergencia.png}
\caption{Error absoluto medio contra $N$. La recta discontinua tiene pendiente
exacta $-1/2$.}
\label{fig:convergencia}
\end{figure}

\begin{table}[H]
\centering\small
\begin{tabular}{llrrr}
\toprule
Integral & Fuente & Pendiente ajustada & Teórica & Desvío \\
\midrule
(a) & MT19937 (propio) & -0.4546 & $-0.5$ & 0.0454 \\
(a) & LCG (propio) & -0.5378 & $-0.5$ & 0.0378 \\
(b) & MT19937 (propio) & -0.5379 & $-0.5$ & 0.0379 \\
(b) & LCG (propio) & -0.5087 & $-0.5$ & 0.0087 \\
\bottomrule
\end{tabular}

\caption{Pendientes ajustadas por mínimos cuadrados sobre los logaritmos.}
\label{tab:pendientes}
\end{table}

Si $|\hat I_N - I| \approx C N^{-1/2}$, entonces
$\log|\hat I_N - I| \approx \log C - \tfrac12 \log N$, de modo que la pendiente
en escala log--log debe ser $-1/2$. Las cuatro pendientes ajustadas quedan
entre $-0.45$ y $-0.54$, con desvíos de a lo sumo $0.046$ respecto al valor
teórico. \textbf{La tasa $O(N^{-1/2})$ derivada en el
Teorema~\ref{teo:varianza} queda confirmada empíricamente.}

Los desvíos no son un defecto del método sino ruido de estimación: el error de
una corrida es una variable aleatoria, y promediar 20 réplicas reduce esa
dispersión pero no la elimina. Con $N$ grande el número de réplicas baja (para
no agotar la memoria), lo que aumenta el ruido justo en el extremo derecho de
la recta, que es donde más pesa el ajuste.

\subsection{Verificación de la cobertura}

El intervalo \eqref{eq:ic} promete atrapar el valor real el $95\,\%$ de las
veces. Eso es comprobable. Se genera una secuencia de $10^7$ uniformes, se
parte en 1000 bloques disjuntos de $N = 10^4$, se construye un intervalo por
bloque y se cuenta cuántos contienen $2/\pi$:
\[
  \textbf{946 de 1000} = 94.6\,\% ,
  \qquad \text{esperado } 950 \pm 6.9 .
\]
El resultado cae dentro de una desviación estándar del valor nominal. Se usan
bloques disjuntos de una misma corrida, y no semillas distintas, porque en un
LCG dos semillas producen secuencias relacionadas de forma determinista; los
bloques disjuntos garantizan que ningún número se reutilice entre réplicas.

\section{Alta dimensión y la maldición de la dimensionalidad}
\label{sec:mc-alta-dimension}

Se estima el volumen de la bola unitaria en $d$ dimensiones,
$V_d = \pi^{d/2}/\Gamma(d/2+1)$, muestreando puntos uniformes en el cubo
$[-1,1]^d$ y contando la fracción que cae dentro. El valor cerrado se usa
únicamente como verdad de referencia.

Como la fracción $p$ es una proporción binomial, $\mathrm{Var}(\hat p) =
p(1-p)/N$, y el intervalo se construye sobre esa varianza multiplicada por el
volumen del cubo, $2^d$.

\begin{table}[H]
\centering\small
\begin{tabular}{rrrrcrrc}
\toprule
$d$ & Aciertos & Esperados & $\hat V_d$ & IC 95\% & $V_d$ exacto & Error rel. & Contiene \\
\midrule
2 & 157\,087 & 157079.633 & 3.141740 & [3.134543, 3.148937] & 3.141593 & 0.00\% & Sí \\
5 & 33\,075 & 32898.681 & 5.292000 & [5.239897, 5.344103] & 5.263789 & 0.54\% & Sí \\
10 & 491 & 498.079 & 2.513920 & [2.291832, 2.736008] & 2.550164 & 1.42\% & Sí \\
20 & 0 & 0.005 & 0.000000 & [0.000000, 15.728640] & 0.025807 & 100.00\% & Sí \\
\bottomrule
\end{tabular}

\caption{Estimación de $V_d$ con $N = 200\,000$ puntos.}
\label{tab:mc-bola}
\end{table}

La columna de aciertos es la que cuenta la historia. En $d=2$ caen dentro
$157\,087$ puntos de $200\,000$; en $d=10$ apenas $491$; en $d=20$,
\textbf{ninguno}. Y no es un error: el número esperado de aciertos en $d=20$ es
$0.005$, así que observar cero es el resultado más probable con diferencia.

\paragraph{El intervalo cuando no hay aciertos.} Con $\hat p = 0$ el error
estándar binomial $\sqrt{\hat p(1-\hat p)/N}$ vale cero y el intervalo
\eqref{eq:ic} colapsa a $[0,0]$, que \emph{no} contiene el valor real
$V_{20} \approx 0.0258$. La aproximación normal exige varios aciertos para ser
válida. En su lugar se usa la \emph{regla de tres}: con cero éxitos en $N$
ensayos, el límite superior del intervalo del $95\,\%$ para la proporción es
aproximadamente $3/N$, lo que da $[0,\;15.73]$. Contiene el valor real, pero es
tan ancho que no informa nada. Esa inutilidad es el diagnóstico.

\subsection{Comparación contra cuadratura determinística}

\begin{table}[H]
\centering\small
\begin{tabular}{rlrrr}
\toprule
$d$ & Puntos & Evaluaciones & $\hat V_d$ rejilla & Error rel. \\
\midrule
2 & $10^{2}$ & 100 & 3.200000 & 1.86% \\
5 & $10^{5}$ & 100\,000 & 4.884480 & 7.21% \\
10 & $10^{10}$ & 10\,000\,000\,000 & inviable & --- \\
20 & $10^{20}$ & 100\,000\,000\,000\,000\,000\,000 & inviable & --- \\
\bottomrule
\end{tabular}

\caption{Cuadratura en rejilla con $m = 10$ subdivisiones por eje.}
\label{tab:mc-rejilla}
\end{table}

Una rejilla de $m=10$ por eje necesita $10^d$ evaluaciones. En $d=2$ son 100 y
en $d=5$ son cien mil, ambas triviales. En $d=10$ son \textbf{diez mil
millones}, y en $d=20$ son $10^{20}$: más evaluaciones que segundos
transcurridos desde el Big Bang. El costo crece exponencialmente con $d$
mientras el refinamiento por eje se queda clavado en 10. \emph{Eso} es la
maldición de la dimensionalidad.

\begin{table}[H]
\centering\small
\begin{tabular}{rrrrr}
\toprule
$d$ & $m$ & Presupuesto & Error rejilla & Error MC \\
\midrule
2 & 447 & 199\,809 & 0.00% & 0.00% \\
5 & 11 & 161\,051 & 2.35% & 0.39% \\
10 & 3 & 59\,049 & 36.68% & 9.56% \\
20 & 2 & 1\,048\,576 & 100.00% & 100.00% \\
\bottomrule
\end{tabular}

\caption{Error a igual presupuesto de evaluaciones: se elige $m$ tal que
$m^d \approx 200\,000$.}
\label{tab:mc-presupuesto}
\end{table}

A igual presupuesto el cruce es nítido. En $d=2$ empatan. En $d=5$ Monte Carlo
ya gana por un factor de seis. En $d=10$ el presupuesto solo alcanza para
$m=3$ subdivisiones por eje, y la rejilla se equivoca por un $36.7\,\%$ contra
el $9.6\,\%$ de Monte Carlo. La rejilla no falla por ser mala regla, sino
porque en alta dimensión el presupuesto no da ni para tres puntos por eje.

\subsection{Por qué el muestreo uniforme se vuelve ineficiente}

\begin{figure}[H]
\centering
\includegraphics[width=0.75\textwidth]{fraccion_volumen.png}
\caption{Fracción del cubo $[-1,1]^d$ ocupada por la bola unitaria, en escala
logarítmica.}
\label{fig:fraccion}
\end{figure}

La fracción $V_d/2^d$ se desploma: $0.785$ en $d=2$, $0.165$ en $d=5$,
$2.5\times10^{-3}$ en $d=10$, $2.5\times10^{-8}$ en $d=20$ y
$2.0\times10^{-14}$ en $d=30$.

El fenómeno geométrico es que \textbf{en alta dimensión, casi todo el volumen
de un cubo está en sus esquinas}. La bola inscrita toca cada cara a distancia
1 del centro, pero los vértices del cubo están a distancia $\sqrt{d}$, que
crece sin límite. El cubo se vuelve, en esencia, $2^d$ púas delgadas alrededor
de una bola comparativamente diminuta.

Esto no contradice la tasa $O(N^{-1/2})$: la tasa sigue siendo $N^{-1/2}$, pero
la \emph{constante} explota. Con $p$ minúsculo, el error relativo del estimador
de proporción es
\[
  \frac{\mathrm{EE}(\hat p)}{p} = \sqrt{\frac{1-p}{Np}} \approx \frac{1}{\sqrt{Np}},
\]
así que mantener un error relativo fijo exige $N \propto 1/p = 2^d/V_d$, que en
$d=20$ significa del orden de $10^8$ puntos para un solo dígito de precisión.

La lección de diseño es que el muestreo uniforme en el cubo es la elección
equivocada: gasta prácticamente todas sus muestras en una región que no aporta.
Un buen estimador en alta dimensión debe concentrar el muestreo donde el
integrando no se anula, que es la idea del muestreo por importancia.

\section{Reducción de varianza}
\label{sec:mc-varianza}

Como el error es $\sigma/\sqrt{N}$, bajarlo admite dos caminos: subir $N$, que
cuesta tiempo lineal, o bajar $\sigma$. Se implementan dos técnicas y se
aplican a \emph{ambas} integrales de la Sección~\ref{sec:mc-integrales}, porque
el contraste entre ellas es el resultado más interesante.

\subsection{Variables antitéticas}

A cada $U$ se le empareja su reflejo $a+b-U$, que tiene la misma distribución.
Con $N/2$ uniformes se producen $N$ evaluaciones, de modo que el presupuesto se
mantiene. Escribiendo $Y = f(U)$ e $Y' = f(a+b-U)$,
\begin{equation}
  \mathrm{Var}\!\left(\frac{Y+Y'}{2}\right)
  = \frac{\mathrm{Var}(Y) + \mathrm{Var}(Y') + 2\,\mathrm{Cov}(Y,Y')}{4}
  = \frac{\mathrm{Var}(Y)\,\bigl(1 + \rho\bigr)}{2},
  \label{eq:antiteticas}
\end{equation}
donde $\rho$ es la correlación entre $Y$ e $Y'$. La técnica ayuda si y solo si
$\rho < 0$, lo que ocurre cuando $f$ es monótona: si $U$ es grande, $a+b-U$ es
pequeño, y una $f$ monótona convierte eso en valores opuestos que se compensan.

\subsection{Variables de control}

Se usa una variable $g(X)$ de media conocida $\mu_g$ y se corrige:
$Y_c = Y - c\,(g(X) - \mu_g)$. Como $\mathbb{E}[g(X)-\mu_g] = 0$, el estimador
sigue siendo insesgado para cualquier $c$. Su varianza es
\[
  \mathrm{Var}(Y_c) = \mathrm{Var}(Y) - 2c\,\mathrm{Cov}(Y,g) + c^2\mathrm{Var}(g),
\]
una parábola en $c$ cuyo mínimo está en
$c^\star = \mathrm{Cov}(Y,g)/\mathrm{Var}(g)$. Sustituyendo,
\begin{equation}
  \mathrm{Var}(Y_{c^\star}) = \mathrm{Var}(Y)\,(1-\rho^2)
  \quad\Longrightarrow\quad
  \text{factor de reducción} = \frac{1}{1-\rho^2},
  \label{eq:control}
\end{equation}
con $\rho$ la correlación entre $f$ y $g$. La técnica vale exactamente lo que
$g$ se parezca a $f$.

\subsection{Resultados}

\begin{table}[H]
\centering\small
\begin{tabular}{llrrrrc}
\toprule
Integral & Técnica & $\hat I$ & Varianza & Factor & \emph{Speedup} & IC 95\% \\
\midrule
(a) & Simple & 0.637478 & 9.4568e-02 & 1.00 & 1.00 & [0.636130, 0.638825] \\
(a) & Antitéticas & 0.637372 & 9.4585e-02 & 0.50 & 0.50 & [0.635466, 0.639279] \\
(a) & Control $g(x)=x$ & 0.637482 & 9.4566e-02 & 1.00 & 1.00 & [0.636135, 0.638830] \\
(a) & Control $g(x)=x(1-x)$ & 0.636617 & 2.9781e-04 & 317.54 & 317.54 & [0.636542, 0.636693] \\
(b) & Simple & 0.478216 & 5.2886e-02 & 1.00 & 1.00 & [0.477208, 0.479224] \\
(b) & Antitéticas & 0.477270 & 7.2851e-05 & 362.97 & 362.97 & [0.477217, 0.477323] \\
(b) & Control $g(x)=x$ & 0.477273 & 4.9911e-04 & 105.96 & 105.96 & [0.477175, 0.477371] \\
(b) & Control $g(x)=x(1-x)$ & 0.477190 & 9.9949e-03 & 5.29 & 5.29 & [0.476752, 0.477628] \\
\bottomrule
\end{tabular}

\caption{Factores de reducción con $N = 200\,000$ evaluaciones en todos los
casos.}
\label{tab:mc-varianza}
\end{table}

\paragraph{¿La técnica siempre ayuda?} No, y las dos integrales lo demuestran
en direcciones opuestas.

\emph{Antitéticas.} En la integral (b), donde $\phi$ es monótona decreciente en
$[0,2]$, el factor es $\mathbf{363\times}$. En la integral (a) el factor es
$\mathbf{0.50\times}$: la técnica \textbf{empeora} el estimador. La razón es
exacta y no estadística: $\sin(\pi(1-x)) = \sin(\pi x)$, de modo que
$Y' = Y$ idénticamente, $\rho = +1$, y \eqref{eq:antiteticas} da
$\mathrm{Var}(Y)$, la misma varianza que Monte Carlo simple, pero habiendo
gastado dos evaluaciones por muestra en lugar de una. La simetría del
integrando respecto al centro del intervalo es precisamente el peor caso.

\emph{Variables de control.} En la integral (a), el control lineal $g(x)=x$ da
$\rho = 0.004$ y un factor de $1.00\times$: no hace nada. Otra vez la simetría
es la culpable, ahora porque $\mathrm{Cov}(\sin(\pi U), U) = 0$ para
$U\sim\mathrm{Unif}(0,1)$. Basta cambiar a un control que comparta la simetría,
$g(x) = x(1-x)$, para obtener $\rho = 0.9984$ y un factor de
$\mathbf{317\times}$. En la integral (b) el orden se invierte: el control
lineal da $106\times$ y el cuadrático solo $5.3\times$.

\paragraph{Contraste con la predicción teórica.} La fórmula \eqref{eq:control}
predice un factor $1/(1-\rho^2)$. Con $\rho = 0.9984$ predice $313$ y se midió
$\mathbf{317}$; con $\rho = -0.9953$ predice $107$ y se midió $\mathbf{106}$.
La coincidencia confirma la derivación.

\paragraph{\emph{Speedup} efectivo.} Como la varianza del promedio es
inversamente proporcional a $N$, reducir la varianza por un factor $k$ equivale
a multiplicar el número de muestras por $k$. El control $g(x)=x(1-x)$ en la
integral (a) rinde como $317 \times 200\,000 \approx 6.3\times 10^{7}$ muestras
de Monte Carlo simple, a costo de una sola evaluación extra por punto.

\section{Cuando el generador arruina el resultado: RANDU}
\label{sec:mc-randu}

Se estima la integral triple
\[
  I = \iiint_{[0,1]^3} \mathbf{1}\bigl[x^2+y^2+z^2 \le 1\bigr]\,dx\,dy\,dz
  = \frac{\pi}{6} \approx 0.52359878,
\]
que depende genuinamente de ternas $(x,y,z)$, usando ternas consecutivas y
disjuntas de cada generador.

\begin{table}[H]
\centering\small
\begin{tabular}{lrcrrc}
\toprule
Fuente & $\hat I$ & IC 95\% & Error & $|$error$|/$EE & Contiene $\pi/6$ \\
\midrule
RANDU & 0.523535 & [0.522843, 0.524227] & -6.38e-05 & 0.2 & Sí \\
Mersenne Twister & 0.523711 & [0.523019, 0.524404] & +1.13e-04 & 0.3 & Sí \\
LCG (Num. Recipes) & 0.522999 & [0.522307, 0.523691] & -6.00e-04 & 1.7 & Sí \\
\bottomrule
\end{tabular}

\caption{Integral triple con $2\,000\,000$ de ternas.}
\label{tab:mc-randu}
\end{table}

\paragraph{Un resultado contrario a la expectativa.} RANDU \textbf{no falla}.
Su error es de $0.2$ errores estándar, mejor incluso que el del LCG de
Numerical Recipes ($1.7$). Los tres intervalos contienen $\pi/6$.

Esto merece explicación en vez de disimulo. Los 15 planos derivados en el
Teorema~\ref{teo:red} tienen normal $\mathbf{h}=(9,-6,1)$ y están separados por
\[
  \frac{1}{\|\mathbf{h}\|} = \frac{1}{\sqrt{81+36+1}} = \frac{1}{\sqrt{118}}
  \approx 0.0921 .
\]
La bola unitaria restringida a $[0,1]^3$ mide del orden de 1 en cada dirección,
así que la atraviesan una docena de planos. Repartir los puntos sobre planos
equiespaciados que cortan la región muchas veces se parece mucho a un muestreo
\emph{estratificado}, y la estratificación no sesga: si acaso, reduce varianza.
El defecto de RANDU no desaparece, simplemente no muerde a esta escala.

\subsection{Dónde sí muerde}

Si la estructura tiene una escala característica de $0.0921$, el fallo debe
aparecer al estimar regiones de ese tamaño o menores. Se repite el
experimento con bolas centradas en $(\tfrac12,\tfrac12,\tfrac12)$ de radio
decreciente.

\begin{table}[H]
\centering\small
\begin{tabular}{rrrrrr}
\toprule
$r$ & $V$ real & RANDU & Error & Mersenne Twister & Error \\
\midrule
0.30 & 1.131e-01 & 1.140e-01 & 0.8% & 1.132e-01 & 0.1% \\
0.15 & 1.414e-02 & 1.480e-02 & 4.7% & 1.422e-02 & 0.6% \\
0.08 & 2.145e-03 & 1.859e-03 & 13.3% & 2.238e-03 & 4.4% \\
0.05 & 5.236e-04 & 7.330e-04 & 40.0% & 5.550e-04 & 6.0% \\
0.03 & 1.131e-04 & 2.750e-04 & 143.2% & 1.260e-04 & 11.4% \\
\bottomrule
\end{tabular}

\caption{Volumen de bolas pequeñas, $10^6$ ternas. El espaciado entre planos
de RANDU es $0.0921$.}
\label{tab:randu-escala}
\end{table}

\begin{figure}[H]
\centering
\includegraphics[width=\textwidth]{randu_convergencia.png}
\caption{Izquierda: los tres generadores convergen para la bola unitaria.
Derecha: al bajar el radio por debajo del espaciado entre planos (línea
punteada), RANDU se despeña mientras Mersenne solo acusa el ruido estadístico
de tener menos aciertos.}
\label{fig:randu-convergencia}
\end{figure}

El error de RANDU pasa de $0.8\,\%$ con radio $0.30$ a $\mathbf{143\,\%}$ con
radio $0.03$, mientras Mersenne Twister solo se degrada de $0.1\,\%$ a
$11.4\,\%$ por el ruido esperable de tener pocos aciertos. El cruce ocurre
alrededor del radio $0.05$, es decir cuando el diámetro de la bola es
comparable al espaciado entre planos.

\paragraph{La explicación geométrica.} Una bola de radio $0.03$ tiene diámetro
$0.06$, menor que la separación $0.0921$ entre planos consecutivos. Según dónde
caiga, o bien ningún plano la atraviesa —y RANDU no puede colocar un solo punto
dentro, estimando volumen cero— o bien un plano la corta por el medio y todos
los puntos de RANDU en esa zona se concentran en un disco, sobrestimando. No
hay término medio, porque \textbf{RANDU no tiene puntos entre los planos}. El
$143\,\%$ de error de la tabla es este segundo caso.

\paragraph{Conexión con la prueba espectral.} Esto es la
Sección~\ref{sec:espectral} vista desde Monte Carlo. Allí se comprobó que
$k_i = 9u_i - 6u_{i+1} + u_{i+2}$ toma exactamente 15 valores enteros, con
desviación $0.00$ respecto al entero más cercano en las $19\,998$ ternas
examinadas. Aquí se mide el precio de esa estructura: el estimador falla cuando
la región de interés es más fina que la retícula. La prueba espectral
\emph{predice} la escala a la que el estimador va a fallar, y la predicción se
cumple.

\paragraph{La moraleja, con un matiz.} Un método de Monte Carlo nunca será
mejor que su fuente de aleatoriedad. Pero el experimento obliga a precisar la
frase: la calidad de un generador no es un número único, sino que depende de la
escala a la que se lo interrogue. RANDU es perfectamente utilizable para
estimar el volumen de la bola unitaria y catastrófico para una bola de radio
$0.03$, con los mismos puntos y en el mismo cubo. Lo peligroso no es que un
generador sea malo: es que sea malo \emph{solo} en el régimen que uno necesita,
mientras pasa todas las pruebas que uno corre.
